\paragraph{Conjugate-coordinate class times a fuzzy factor: fibre-point forms, the level-one two-point function with the tree form, the equal-weight elimination of the fibre harmonics and the static response (conditional; bounded obstruction).}
This paragraph supplements the preceding one, in which the elimination with the tree-level form, the trace polarisation and a Newtonian limit were not computed. Keep its setting and notation, and take \(q=(q_0,0,0,k)\) at \(y^i=0\), so that \(q^2=k^2-q_0^2\), \(a=m^2q_0^2+\ell^2k^2\), \(b=m^2(q_0^2-k^2)\), \(\sinh^2\tau=m^2/\ell^2\) and \(\vartheta=1+\ell^2/m^2\). Below \(\alpha\) is any function of the fibre point with values in the frame directions, modulo the gauge directions \(\gamma(n)q\), so that \(q\cdot\alpha=0\) can be taken. All forms are for a real wave averaged over its phase.

\emph{Fibre-point forms.} At a fibre point, exactly, \(X_n=m^{-4}q^2(v\cdot v)(\alpha\cdot\alpha)\) and the scalar density is \(\mathrm{SC}_n=-\tfrac1{24}m^{-4}q^2[(v\cdot v)(\alpha\cdot\alpha)+2(\alpha\cdot v)^2]\). With \(h_n=m(\alpha\otimes v+v\otimes\alpha)\) and \(t_n=\eta_{\mu\nu}h_n^{\mu\nu}\) these are \(\tfrac12m^{-6}q^2(h_n\cdot h_n-\tfrac12t_n^2)\) and \(-\tfrac1{48}m^{-6}q^2(h_n\cdot h_n+\tfrac12t_n^2)\). Since \(q\cdot\alpha=q\cdot v=0\), the form \(L_2\) of the preceding paragraphs vanishes on all nine amplitudes \(D_{ij}\) with \(\gamma_i=0\). For \(\alpha=n^ic_i\), \(c_i=D_{ij}s_j\), the fibre average gives, for every \(D\), \(\mathrm{IX}_2=\tfrac1{15}m^{-4}q^2\sum_{ij}(\operatorname{tr}G\,\delta_{ij}+2G_{ij})\,c_i\cdot c_j\) and \(\mathrm{SC}_2=-\tfrac1{360}m^{-4}q^2[\operatorname{tr}G\operatorname{tr}P+2\operatorname{tr}GP+2(\operatorname{tr}A)^2+2\operatorname{tr}AA^{T}+2\operatorname{tr}A^2]\), with \(P=DGD^{T}\) and \(A=DG\) (exact rational arithmetic at nine cases, two with \(q\cdot y=0\)). These forms are polynomial in the amplitude \(c\). The term \(1/x\) of the preceding paragraph comes from dividing by \(\bar h\cdot\bar h\), which on helicity \(\pm2\) is \(\tfrac49m^2a\,c_i\cdot c_i\), so it needs no eliminated tower to explain it. Exactly, \(\operatorname{tr}G\,\delta_{ij}+2G_{ij}=5m^2[\delta_{ij}(q_0^2-\tfrac15k^2)-\tfrac25k_ik_j]+\ell^2k^2(4\delta_{ij}-2\hat q_i\hat q_j)\), so \(\mathrm{IX}_2\) has, up to the terms suppressed by \(1/\sinh^2\tau\), the bracket of Equation~(4.41) of~\cite{Steinacker2026}, including its transversal factor \(q_0^2-\tfrac15k^2\) and longitudinal factor \(q_0^2-\tfrac35k^2\); that work has the torsion term only, so the comparison is of \(\mathrm{IX}_2\) and of the structure, and the overall normalisations are not compared. With \(n\cdot\hat q=\cos\theta\) one has \(v\cdot v=a\sin^2\theta+b\cos^2\theta\), and the numbers \(\tfrac13\), \(\tfrac15\), \(\tfrac35\) are the averages of \(\cos^2\theta\) with the weights \(1\), \(n_\perp^2\) and \(n_\parallel^2\): they are moments of one fibre function, up to normalisation the symbol of the operator \(\tilde\Box\) of that work, taken on different fibre profiles.

\emph{Helicity \(\pm2\) at level one with the tree form.} Take \(S_E=\tfrac1{4g^2}\operatorname{Tr}F_{ab}F^{ab}\) with \(\operatorname{Tr}f=\int\mathrm d^4y\,\rho_M\langle f\rangle\), \(\langle\cdot\rangle\) the fibre average, so that the leading-symbol tree form is \(\tfrac{\rho_M}{4g^2}m^2(q^2\langle\alpha\cdot\alpha\rangle-\langle(q\cdot\alpha)^2\rangle)\); the coupling and \(\rho_M\) are supplied. With the weights of the preceding paragraph the induced coefficient relative to \(\mathrm{EH}(\bar h)\) on helicity \(\pm2\) is \(c_T^{\mathrm{TT}}+\kappa/x\), \(\kappa=\tfrac1{20}(V_1-T_1)+\tfrac3{20}T_2\ge\tfrac1{20}(V_1+2T_1)\). In the formal continuation \(S_L=-S_E\) with the Lorentzian \(\eta\), for \(c\cdot c=\sum_ic_i\cdot c_i\),
\[
S_L=c\cdot c\;Z\,(q_0^2-k^2)(q_0^2-w^2),\quad
Z=\frac{c_T^{\mathrm{TT}}+\kappa}{576\pi^2m^2},\quad
w^2=v^2k^2-\mu^2,\quad
v^2=\frac{\kappa-c_T^{\mathrm{TT}}\ell^2/m^2}{c_T^{\mathrm{TT}}+\kappa},\quad
\mu^2=\frac{48\pi^2\rho_Mm^4}{g^2(c_T^{\mathrm{TT}}+\kappa)},
\]
exactly for these forms, and \(1-v^2=c_T^{\mathrm{TT}}\vartheta/(c_T^{\mathrm{TT}}+\kappa)>0\). For \(V_i=T_i=T\), \(v^2=(m^2-4\ell^2)/5m^2\); the value \(\tfrac15\) at large \(\tau\) needs these weights, and in general the limit is \(\kappa/(c_T^{\mathrm{TT}}+\kappa)\). The factorisation and the sign of the tree term agree with Equation~(4.44) of~\cite{Steinacker2026}, whose extra factor has the mass term of Equation~(4.45), not compared in size with \(\mu^2\). At fixed \(k\) the coefficient of \(q_0^2-k^2\) at the massless root is \(Z(k^2-w^2)>0\) and that of \(q_0^2-w^2\) at the extra root is \(-Z(k^2-w^2)<0\), with \(k^2-w^2=(1-v^2)k^2+\mu^2\). The propagator of \(\bar h\) has the extra factor \(\tfrac49m^2a\), with \(a=m^2\vartheta k^2\) at the massless root, which gives the residue \(256\pi^2m^6/c_T^{\mathrm{TT}}\) for \(k^2\gg\mu^2\), independent of \(\kappa\), and \(a=m^2(w^2+k^2/\sinh^2\tau)>0\) at the extra root for \(w^2\ge0\). So within these hypotheses the massless root has the healthy sign and the extra root has the opposite sign in both variables. It has real frequency only for \(v^2k^2\ge\mu^2\) and is tachyonic otherwise. The one-loop term is \(-1\) times the tree term at the extra root for every \(k\), and \((1-v^2)k^2/\mu^2\) times it on the null cone, so the root is not inside an expansion in which the one-loop term is small compared with the tree term, while the Einstein regime on the null cone needs \((1-v^2)k^2\gg\mu^2\). Whether it is a propagating state therefore depends on the control of higher loops, for example by a large number of modes of \(K\), which is not established here. Adding \(-Z(V_{00}+V_{01}q_0^2)\), with the signs of Equation~(4.46) of~\cite{Steinacker2026}, does not change the opposite signs while the roots are real and distinct, since they are the signs of the derivative of a quadratic in \(q_0^2\) at its two roots; a complex pair needs \((k^2-w^2)^2+2V_{01}(k^2+w^2)+V_{01}^2+4V_{00}<0\), and to first order the root near \(k^2\) moves by \((V_{00}+V_{01}k^2)/(k^2-w^2)\). These terms are not computed here. That work states that the extra modes raise no issue of ghost-like signatures because they arise only after quantization; this is neither derived nor contradicted here.

\emph{Equal weights and all fibre harmonics.} Suppose all fibre transitions have the same proper-time weight, the equal-weight trace of the preceding paragraphs, which holds when \(M\) exceeds the fibre levels that enter; then \(V_i=V\), \(T_i=T\), with \(T>0\) for a non-zero flux and \(V>0\) and finiteness assumed. The induced form is then a multiplication operator on the fibre, so the elimination of the harmonics, which the preceding paragraphs left open, is pointwise. With \(C=q^2A/32\pi^2m^4\), \(A=\tfrac1{24}(V+2T)\), \(A+B=\tfrac18V\), the form at \(n\) is \(C[(v\cdot v)(\alpha\cdot\alpha)+\tfrac BA(\alpha\cdot v)^2]\) plus the tree term \(\tfrac{\rho_M}{4g^2}m^2q^2\alpha\cdot\alpha\). The component of \(\alpha\) orthogonal to \(v\) has one coefficient at every fibre point, \(4A=\tfrac16W_R-4W_X\) with \(W_R=V-T\), \(W_X=-T/8\), in the sense that its form is \(-\tfrac1{32\pi^2}\cdot4A\,\mathrm{EH}[h_n]\); for \(V=T\) the whole form at \(n\) is \(-\tfrac1{32\pi^2}\cdot\tfrac T2\) times the form \(-\tfrac18m^{-6}q^2(h_n\cdot h_n-\tfrac12t_n^2)\), which is the Einstein--Hilbert form of \(\hat h_n=h_n-\tfrac12\eta t_n\). For \(V=T\) the level-one extra root is a Rayleigh quotient of a continuum of roots \(q_0^2=k^2(\cos^2\theta-\sin^2\theta/\sinh^2\tau)-\varpi/m^2\), \(\varpi=8\pi^2\rho_Mm^6/g^2A\), over all \(\theta\). On polarisations with \(\alpha\cdot\alpha>0\) their sign is the one found above, and on \(\alpha\cdot\alpha<0\) the opposite. So the obstruction to a two-point function without a negative residue persists at equal weights, for the same reasons and with the same caveat.

Let a symmetric \(X\) couple through \(X_{\mu\nu}\bar h^{\mu\nu}=\langle2m\,\alpha\cdot Xv\rangle\). Eliminating \(\alpha\) gives \(-32\pi^2m^6I[X]/q^2A\) with, for \(\varrho=A/(A+B)=(V+2T)/3V\),
\[
I[X]=\Bigl\langle\frac{Xv\cdot Xv}{v\cdot v}\Bigr\rangle+(\varrho-1)\Bigl\langle\frac{(v\cdot Xv)^2}{(v\cdot v)^2}\Bigr\rangle,
\]
a principal or finite-part value where \(v\cdot v\) vanishes; this is the real part of the prescription \(q_0^2\to q_0^2+\mathrm i\epsilon\), whose imaginary part is not treated. The form \(c\,\mathrm{EH}\) with one coefficient \(c\), on the variable \(\hat h=\bar h-\tfrac12\eta\operatorname{tr}\bar h\), would give \(-64\pi^2m^6(X\cdot X-(\operatorname{tr}_3X)^2)/cq^2\), with \(\operatorname{tr}_3\) the trace on the three-space orthogonal to \(q\). On the null cone with \(X\) transverse and traceless, \(I=\tfrac14(1+\varrho)X\cdot X\) exactly, so the coefficient is \(c^{\mathrm{TT}}=8A/(1+\varrho)=V(V+2T)/2(2V+T)\), which is \(T/2\) at \(V=T\) against the level-one value \((2V+T)/5\), which is larger for all \(V,T>0\). A truncation to fibre orders \(J\le N\) reproduces the level-one value at \(N=1\) and tends to this one (numerical, up to \(N=25\)).

\emph{Static response.} Matter coupling is not derived in this ledger, so two couplings are recorded for dust \(T_{\mu\nu}=\varrho_0u_\mu u_\nu\), \(u\cdot u=-1\), \(q\cdot u=0\). (S1) The source couples to \(\hat h\) as \(T_{\mu\nu}\hat h^{\mu\nu}\), the form of the matter variation in Equation~(3.9) of~\cite{Steinacker2026}, so that \(X=T-\tfrac12\eta\operatorname{tr}T\); it then satisfies \(X\eta^{-1}X=\tfrac14\varrho_0^2\eta\), hence \(Xv\cdot Xv=\tfrac14\varrho_0^2v\cdot v\). (S2) The source couples to \(\bar h\) with \(X=T\), the Poisson-limit source of a mode with action \(\int\rho_M(\gamma^{\mu\nu}\partial_\mu\phi\partial_\nu\phi+M^2\phi^2)\) and a rigid mass from \(K\). Put \(c_0=1/\cosh\tau\) and \(g(\tau)=\vartheta^{-1}(1-c_0\operatorname{arctanh}c_0)\), the principal value of \(\langle(v^0)^2/(-v\cdot v)\rangle\) in the rest frame of \(y\) at \(q\cdot y=0\); it tends to \(1\) with \(1-g\simeq2/\sinh^2\tau\).

At level one, for \(q\cdot y=0\) and \(D=\mathrm{diag}(d_\perp,d_\perp,d_q)\), the block with helicity \(0\) and the trace polarisation, the response with the weights of the preceding paragraph has, as \(\ell/m\to0\) at fixed weights, a ratio to the same source with the coefficient \(c_T^{\mathrm{TT}}\) of \(10c_T^{\mathrm{TT}}/(5V_0+4(V_1-T_1)+4T_2)>0\) for (S1) and four times this for (S2). The helicity-\(0\) eigenvalue relative to \(\mathrm{EH}(\bar h)\) does tend to \(-\sinh^2\tau[\tfrac1{60}(V_1-T_1)+\tfrac1{20}T_2]<0\), but the source does not act on that polarisation alone: tracelessness forces \(d_\perp=m^2d_q/2\ell^2\), whereas at large \(\tau\) the source acts on \(d_q\), whose diagonal entry in the block is positive at large \(\tau\). So at late times the sign of the static response is that of \(c_T^{\mathrm{TT}}\). It is not so at every \(\tau\): for \(V_i=T_i\) and (S1) the ratio is \(2(\sinh^4\tau-10\sinh^2\tau+4)/(\sinh^2\tau-4)(3\sinh^2\tau-2)\), with poles at \(\sinh^2\tau=\tfrac23,4\), zeros at \(\sinh^2\tau=5\pm\sqrt{21}\), negative for \(4<\sinh^2\tau<5+\sqrt{21}\), tending to \(\tfrac23\); and for (S2) it is \(8\sinh^2\tau/(3\sinh^2\tau-2)\), tending to \(\tfrac83\).

With all harmonics and equal weights the level-one defect is removed for (S1) at \(V=T\): then \(\varrho=1\) and \(I=\tfrac14\varrho_0^2\) for every \(q\) with \(q\cdot u=0\), every base point and every \(\tau\), without a principal value, which is the value of the one-coefficient form at \(c=T/2=c^{\mathrm{TT}}\) since \(X\cdot X-(\operatorname{tr}_3X)^2=\tfrac12\varrho_0^2\). So at \(V=T\), and with that coupling, the static dust response and the transverse-traceless pole have one coefficient, with the attractive sign. Solving for \(\alpha\) at rest gives \(\hat h^{00}\) independent of \(\tau\) and \(\hat h^{11}/\hat h^{00}=g(\tau)\), which is the ratio of the two potentials once the part along \(\hat q\hat q\) is removed as a gauge term, and it tends to \(1\) as \(\tau\to\infty\). For \(V\neq T\) and (S1), \(I=\tfrac14\varrho_0^2[1+(\varrho-1)K(\tau)]\) with \(K\) a closed function tending to \(1\), so at large \(\tau\) the coefficient is \(V/2\) and the response is \((V+2T)/(2V+T)\) times that of \(c^{\mathrm{TT}}\); that ratio of potentials still tends to \(1\) (numerical at two values of \(\varrho\)). For (S2) at \(V=T\), \(I=\varrho_0^2g(\tau)\): the response is \(4g(\tau)\) times that of the coefficient \(c^{\mathrm{TT}}\), and \(\bar h^{ab}=0\) on the two directions orthogonal to \(q\) and \(u\), so the second potential vanishes and that ratio is \(0\). So within these hypotheses the value \(1\) of that ratio depends on the coupling of massive modes to the scale of \(K\): (S1) is the source of the same action with \(M^2\propto\rho/\sqrt{|g|}\), \(\rho\) the density of the preceding paragraph, which is not derived. Beyond dust the agreement fails: for \(T=\varrho_0uu+p(e_1e_1+e_2e_2)\) in (S1) at \(V=T\), \(I=\tfrac14\varrho_0^2+g(\tau)(\varrho_0p+p^2)\), against \(\tfrac14\varrho_0^2+\varrho_0p\), and on helicity \(\pm2\) at \(V=T\), away from the null cone, \(I\) carries the factor \(\tfrac{x}{x-1}\bigl[1-\operatorname{arctanh}\sqrt{(x-1)/x}\,/\sqrt{x(x-1)}\bigr]\) for \(x>1\), and \(1-g\) at \(q\cdot y=0\), so a single coefficient for all \(q\) and all sources is not obtained. With the tree term at \(V=T\), (S1) and \(q\cdot y=0\), the factor is \(\langle v\cdot v/(v\cdot v+\varpi)\rangle\) as a principal value, which is the fibre-resolved form of \(\tilde\Box/(\tilde\Box-m^2)\) of Equation~(3.25) of~\cite{Steinacker2026}; for \(\varpi\ll m^2k^2\) at large \(\tau\) it is \(1-\varpi/m^2k^2\) to first order, against \(1+3\varpi/m^2k^2\) for the fibre average of \(\tilde\Box\) with which that work reads its static halo, and it tends to zero as \(k\to0\) as there. This difference of sign is conditional on these hypotheses and on dropping the fibre Laplacian, which by a scale estimate requires \(\varpi\gg1\); it may be an artefact of the frozen symbol or of the equal-weight trace, and that work itself states that the higher-spin parts were ignored in that derivation.

These statements are for the hypotheses of the preceding paragraph, zeroth order in the background torsion, the leading symbol for the loop and for the tree form, which drops non-derivative terms and the fibre Laplacian and on the null cone holds up to a shift of relative order \(\vartheta/m^2k^2\), the formal continuation with the Lorentzian \(\eta\), and, for the elimination of all harmonics, a weight independent of the fibre orders of the external amplitudes, which is an assumption outside \(M\) above the levels that enter. The forms are exact as stated; the eliminations are exact at these hypotheses apart from the principal values and the numerical checks stated. They do not compute the fibre Laplacian, the cross term, higher loops or \(V_{00},V_{01}\), do not derive a matter coupling, the scale of \(K\) or the condition \(V=T\), and do not construct \(K\). They do not compare the modified Einstein equations of~\cite{KumarSteinacker2023,Steinacker2026} beyond the equations named, do not identify \(\bar h\), \(\hat h\) or \(h_n\) with the emergent metric of Chapter~\ref{ch:vierbein-emergence}, and do not bear on Assumption~\ref{ass:lorentzian}, so no value of \(G_{\mathrm{eff}}\), no bound of Cassini type and no statement on the presence or absence of a propagating ghost for the emergent metric follow. Problems~C, F and~N remain open and PROGRAMME-CONDITIONAL.
